% ════════════════════════════════════════════
% 分析目标（02.分析目标.tex）
% ════════════════════════════════════════════
\section{分析目标}
\subsection{总体目标}

本报告的总目标是：在同一份帖子数据上识别曝光量的主要影响因素，量化付费推广与内容选择带来的曝光差异，并把统计结果转化为可执行的发帖与投放建议。目标不止于拟合一个高拟合优度的模型，而是要回答三个决策问题：哪些变量在发帖前可以调整，哪些变量只是账号成长的伴随标签，调整后预计能带来多少曝光增量。

\subsection{具体目标}

围绕总体目标，报告设定四项具体目标。第一，完成数据理解与分布诊断，判定变量单位、数据粒度、截面或时序结构及因果角色，量化曝光量的重尾程度与粉丝规模、月份之间的共线结构，为变量取舍划定边界。第二，完成数据清洗与特征整理，处理缺失、核查异常、统一编码并对曝光量做合适变换，形成可回归、可追溯的分析数据，并通过缺失处理对照检验清洗方案的敏感性。第三，建立曝光量对粉丝规模、帖子类型、内容主题、月份、星期、时段与付费标记的回归模型，给出系数、显著性、拟合优度与区间估计，完成共线性、残差、整体显著性与样本外预测能力检验。第四，基于最终模型测算付费推广的额外曝光、视频相对图片的曝光倍数以及内容主题的显著性与两两差异，并在多种模型设定下复核结论的稳健性，进一步检验付费效应是否随内容主题变化。

\subsection{分析范围与边界}

本报告使用同一主页的 500 条历史帖子，结论适用于该主页的观测期，不宜直接外推到其他品牌、其他平台或未来时段。数据为截面观测，缺少帖子级的随机实验，付费效应通过控制协变量后的反事实比较得到，应理解为条件平均差异而非严格因果效应。曝光量只反映分发结果，不包含点击、互动与转化，因而相关建议以曝光为目标，互动与转化仍需后续数据验证。
